% Revised 23 September 2026 from the supplied methodology draft.
% Retrieval and reconciliation are reported as completed; corpus-wide screening
% and synthesis results remain unfinished. E7 follows executed protocol v2.0.0.
% Removing E7 as a manually scored gate was discussed but is not reported here
% as an implemented change. Update this account if a later amendment is executed.

\section{Review Methodology}
\label{sec:methodology}

The review examines how computational assistance participates in life-science visual analysis, which tasks it supports, and what evidence establishes its consequences for scientific work. Its scope follows the relationship among the data, the visual representation, the computational contribution, and the scientist's role. The use of an AI method or an immersive display alone does not establish that relationship.

The protocol separates scientific eligibility, administrative verification, and classification for synthesis. Scientific screening identifies candidate systems; administrative verification establishes publication and access requirements; classification describes the tasks, assistance mechanisms, and environments of qualifying systems. The locus of analytic agency guides interpretation of computational interventions and human responsibilities rather than serving as an additional eligibility score. The original coding rubric was dated 30 August 2026. A subsequent versioned amendment, implemented as title-and-abstract screening protocol v2.0.0, deferred administrative checks and separated screening outcomes from follow-up actions.

\subsection{Review Scope and Eligibility}
\label{sec:methodology-scope}

The screening unit is a bibliographic paper record. Where a paper describes multiple systems, at least one identified system must jointly satisfy the scientific criteria; evidence about unrelated systems cannot be combined to establish eligibility. Systems and their mechanisms provide units for subsequent comparison, but individual mechanisms are not counted as additional papers.

\begin{description}
\item[E1: Life-science application.] The reported system is applied to analysis, understanding, monitoring, decision-making, discovery, or scientific practice in a life-science or closely related domain, including biomedical and clinical science, health, neuroscience, ecology, marine science, biology, and bioinformatics. Clinical and surgical applications follow the same rules as other applications. Training, teleoperation, procedural guidance, or presentation alone does not establish a qualifying analytic workflow.

\item[E2: Relational, derived-structure, or multiscale relevance.] The system analyzes explicit networks or relationships, or relationships derived from spatial, temporal, multivariate, image-derived, lineage, similarity, or other multiscale life-science data. A node--link representation is not required. Derived similarities and spatial relationships may qualify when they have an analytic role; the presence of image data or a network term alone does not establish that role.

\item[E3: Visual-analytics component.] Humans use an interactive or analytically meaningful visual representation to inspect, understand, validate, steer, compare, interpret, or act on data or computational results. Visual material used solely for illustration, presentation, communication, or artistic output does not qualify. An unclear description of the visual representation's analytic role leaves the criterion unresolved.

\item[E4: Substantive computational assistance.] A computational mechanism contributes nontrivially within the visual-analytics workflow. Examples include clustering, dimensionality reduction, graph computation, layout, aggregation, annotation, ranking, prediction, natural-language analytic operations, and workflow orchestration. Ordinary rendering, pan, zoom, literal manual selection, direct parameter entry, and ordinary database lookup alone do not qualify. Labels such as ``AI,'' ``adaptive,'' or ``agentic'' require supporting evidence of computational behavior and its analytic role.

\item[E5: Human analytic relationship.] Humans meaningfully inspect, direct, validate, interpret, collaborate with, supervise, correct, or make decisions from the assisted computational process. Fully offline or automated computation without a meaningful human analytic relationship to its results is excluded. Manual control over every operation is not required.
\end{description}

Administrative scope, retained as E6, requires a journal, conference, workshop, or preprint research paper with English-language full text available for formal assessment. There is no lower publication-date cutoff; the documented retrieval cutoff is 20 September 2026. A survey or conceptual paper qualifies only when it independently reports or analyzes a system satisfying E1--E5. Publication community is not an additional restriction. E6 is not scored during title-and-abstract screening: it remains explicitly unassessed until administrative verification. An English abstract does not establish English full-text availability, and a failed download does not establish a language-based exclusion.

Protocol v2.0.0 retains E7 as an evidence-sufficiency judgment over E1--E5, rather than an additional scientific characteristic of a system. Each scientific criterion receives Yes, No, or Uncertain. A defensible No in E1--E5 produces Excluded, even when other responses are uncertain. Yes for all five criteria with sufficient evidence produces Include, meaning retained for further assessment rather than final inclusion. Otherwise, a record without a scientific No remains Uncertain. E7=No or Uncertain cannot alone establish scientific exclusion. Incomplete or invalid responses are identified separately by return validation.

The same scientific rules apply across retrieval sources, assistance modes, tasks, and display environments. Prior-survey membership, historical relevance scores, and synthesis priority do not confer eligibility. General frameworks and other nonqualifying works may inform the discussion as background while remaining outside the formal eligible-corpus counts.

\subsection{Sources, Search Strategy, and Retrieval}
\label{sec:methodology-search}

Candidate identification combined PubMed, Europe PMC, Semantic Scholar, IEEE Xplore, and ACM Digital Library retrieval with a local arXiv metadata-snapshot search and prior-survey materials. These routes address differences in disciplinary coverage and terminology: a system may be described through its scientific application, computational method, interface, or display environment. The sources overlap, and retrieval through several routes does not imply several independent studies.

Five overlapping query families covered: (QF01) relational visualization; (QF02) assisted visual analysis; (QF03) interactive systems; (QF04) nondesktop environments; and (QF05) conversational analysis. Each combined life-science terms with its relevant visualization, system, relationship, assistance, or environment expressions. Assistance-neutral routes were retained to identify systems whose computational contribution is not advertised through AI terminology. These families are retrieval strategies, not mutually exclusive paper categories or substitutes for eligibility assessment.

Source-specific acquisition used paginated APIs and ACM search exports, with retained query definitions, source records, and checkpoints. Provider syntax, searched fields, phrase handling, and indexing differ; shared conceptual queries therefore do not guarantee equivalent result sets. Full executed expressions and source-specific acquisition details are to be reported in the supplementary search record.

Repeated arXiv API failures produced no accepted API records. The approved replacement route searched version 303 of the arXiv metadata snapshot under an explicitly defined local matching policy. A single pass over 3,164,528 metadata records yielded 958 distinct matched records and 1,333 record/query-family memberships. Matching used Unicode NFKC normalization, case folding, whole alphanumeric tokens, contiguous within-field phrases, and explicit Boolean grouping, without stemming or fuzzy matching. Equivalence to arXiv API search behavior was not established. The snapshot's observed metadata-update dates extended through 11 September 2026; this observation does not prove exhaustive coverage through that date or through the overall 20 September retrieval cutoff.

Prior-survey identification used materials associated with the surveys by Ehlers and colleagues, Joos and colleagues, and Fonnet and Pri\'e. These materials were treated according to their provenance: the Ehlers-associated working workbook supplied 816 row occurrences, the Joos-associated companion material supplied 146 occurrences, and reconstruction from the Fonnet and Pri\'e bibliography supplied 207 occurrences. The latter included 132 supported survey-member labels, 43 unconfirmed entries, and 32 background references. Working-ledger inclusion was not treated as confirmed survey membership, and bibliography reconstruction was not represented as an exact reconstruction of a surveyed corpus. All candidates remained subject to this review's own eligibility criteria. These imports do not establish that recursive backward or forward citation searching was performed.

% Insert the three survey citations using the manuscript's verified BibTeX keys.
% Supplement: executed query strings/fields, dates, filters, source completion
% evidence, seed-artifact references, and snapshot version/matching specification.

\subsection{Occurrence Preservation and Deduplication}
\label{sec:methodology-deduplication}

The registered global merge preserved 187,446 retrieved or imported occurrences and grouped them into 140,959 canonical bibliographic records before eligibility screening. Each occurrence retains its source and retrieval or import context, including query membership and available identifiers or artifact locations. These counts describe the identification pool, not an eligible corpus or a count of independent studies.

Generic canonicalization used normalized DOI keys when available and exact normalized-title keys otherwise. It did not perform fuzzy title matching, and a DOI-bearing record did not automatically merge with a DOI-missing record solely because their titles agreed. Approved identity decisions supplemented this grouping: 424 arXiv-identifier adjudications reduced the generic total from 141,383 to 140,959 canonical records. Compatible title-group occurrences were preserved when those groups joined the resolved targets. Superseded decisions remained in the history rather than being overwritten.

Thirty prior-survey identity groups covering 71 occurrences were subsequently confirmed through a separately registered adjudication overlay. This confirmation preserved their existing canonical membership, source metadata, and historical unresolved markers. It neither confirmed survey membership nor established scientific eligibility. Effective identity status requires the base dataset and its validated overlay to be interpreted together.

Related-version evidence was retained separately: 50 links covering 45 candidate records did not themselves trigger duplicate consolidation. Consequently, canonical records may include related preprints and published reports. Later synthesis must distinguish paper counts from systems and independent evaluations. The registered merge, identity overlays, input inventories, and normalization implementation provide the derivation of the identification counts; implementation-specific details belong in the reproducibility materials rather than the eligibility rules.

\subsection{Screening, Follow-up, and Human Decisions}
\label{sec:methodology-screening}

Title-and-abstract screening uses the evidence supplied for each record. Explicit title evidence may support an individual criterion, but a promising title does not establish every criterion. Missing information is represented as uncertainty in the affected judgment rather than as a scientific failure. Criterion responses, aggregate screening outcomes, and next actions are recorded separately.

A clear scientific exclusion requires no further acquisition solely because unrelated administrative or scientific information is unknown. Non-excluded records without abstracts enter metadata recovery before human full-text assessment. Conflicting title/abstract evidence is held for reconciliation. Other retained candidates proceed to targeted second review or full-text assessment as appropriate to the unresolved question. An Include outcome remains provisional until scientific and administrative eligibility are established from adequate evidence.

Metadata recovery retrieves existing bibliographic evidence rather than generating substitute abstracts. Recovered material must match the paper's identity and retain its source and acquisition provenance. Original inputs and judgments are preserved; review after enrichment is a distinct evidence version. Missing abstracts, inaccessible full text, parsing failures, and unresolved identity conflicts are not themselves proof of scientific ineligibility. Records left unresolved at the review stopping point must be reported separately from established exclusions. The procedure does not assume that every incomplete record can be recovered or receive human full-text assessment.

Where LLM screening is used, outputs are recorded as proposals with the reviewed input, protocol, model and prompt versions, technical status, criterion responses, and rationales. Model output does not by itself establish final corpus membership. Human corrections and any adjudication retain the earlier judgments and identify the actor, evidence, and reason for the change. Independent double screening is reported only for records that actually receive independent judgments; repeated or model-assisted assessments are not counted as independent reviewers.

The completed development pilot used a frozen sample of 100 canonical records: 70 uniformly sampled records and 30 challenge records selected using predefined metadata and provenance rules. The challenge sample comprised six records each from prior-survey routes, missing abstracts, sparse abstracts, uncommon modality/assistance cues, and document-type boundary cues, with a fixed overlap priority. Ten additional records formed a disjoint calibration set. Two optional, nonoverlapping 25-record secondary assignments were drawn from the same 100 evaluation records; assignments alone do not establish completed independent review. Known related-version links were used to prevent calibration/evaluation splitting without changing bibliographic identities.

The archived primary v2.0.0 return contained 5 Include, 41 Uncertain, and 54 Excluded judgments. One random-sample record was subsequently held because its title and abstract described different papers, leaving 99 records available for descriptive comparison pending reconciliation and re-review. Actionable routes were 31 metadata-recovery cases, 15 full-text-assessment cases, 53 records requiring no further action, and one evidence-reconciliation case. These are pilot judgments and routes, not final eligibility outcomes or corpus-wide workload estimates. Random and challenge results must be analyzed separately. Exposure to prior model outputs and protocol-development judgments prevents this pilot from being described as an independent blinded reference or as evidence of corpus-wide screening accuracy.

\noindent\textbf{Completion required before submission:} Report the executed corpus-wide screening procedure, model and prompt settings if used, human review coverage, recovery and assessment stopping rules, and final dispositions. Any further change to E7 or outcome aggregation must be versioned and described as implemented rather than inferred from discussion. Recovery of missing abstracts and reconciliation of the pilot conflict were still in progress at this revision.

\subsection{Multilabel Coding and Synthesis}
\label{sec:methodology-synthesis}

Coding describes qualifying systems through the task, assistance, and modality perspectives introduced in Section 3. Labels are assessed independently as Present, Absent, or Uncertain. Systems may support several tasks, combine assistance modes, and provide coordinated views in multiple environments; neither a primary category nor a tie-break between overlapping labels is required. A general claim of exploration or intelligence does not establish every task or assistance mode.

The five task categories are Navigation and Multiscale Orientation; Selection, Filtering, and Precision Interaction; Comparison and Differentiation; Sensemaking and Hypothesis Development; and Coordination and Collaborative Reasoning. They identify the analyst's intended action or accomplishment. Provenance and process awareness are treated as workflow or evaluation properties unless reported behavior establishes how they support a specific task.

Assistance coding identifies the computational contribution. Algorithmic assistance operates on data, structure, representations, or analytic candidates. Adaptive assistance tailors views, interaction, guidance, or actions using behavior, context, data characteristics, or inferred intent. Conversational assistance interprets language or dialogue within analytic interaction. Immersive assistance uses spatialization, embodiment, navigation, tracked context, or multisensory cues as part of a computational assistance mechanism. These modes can overlap. A headset or direct controller manipulation alone does not establish Immersive assistance; speech used solely as a button-equivalent command does not establish Conversational assistance.

Modality identifies the environment in which analytic visualization is presented and used: Desktop/Planar, Large Display, VR, AR/MR, or CAVE. It is separate from assistance and device form factor. A rotatable 3D scene on a conventional monitor remains Desktop/Planar, and a controller alone does not add a modality. Multiple labels require evidence of actual analytic use in each environment. Desktop/Planar is the manuscript-facing name for the rubric's historical Desktop 2D storage label.

The coding scheme links assignments to supporting passages and source locations. Assistance annotations identify the system and mechanism, relevant input or signal, computational behavior, output or action, human role, and analytic consequence where the evidence supports these details. A reported capability is distinguished from evidence that it improves analysis. Comparisons examine mechanisms within their system and task context; paper-level label unions must not imply that every mechanism supports every task or environment.

Core, Supporting, and Contextual priorities guide the depth and placement of synthesis after eligibility assessment. They are not study-quality scores or an additional inclusion filter. Noneligible background references remain outside the formal corpus. The locus of analytic agency guides comparison of who proposes, scopes, executes, and evaluates an operation, and how scientists inspect, revise, or reverse computational contributions. These responsibilities are interpreted from documented workflows rather than converted into another categorical score.

\noindent\textbf{Completion required before submission:} Identify the coding actors, material reviewed, actual checking and disagreement-resolution procedures, and final paper/system/mechanism denominators. The coding framework above specifies the method; completed extraction and synthesis validation are not asserted here.

\subsection{Review Flow, Reproducibility, and Limitations}
\label{sec:methodology-limitations}

The review flow distinguishes identification occurrences, canonical bibliographic records, title-and-abstract screening, metadata recovery, full-text assessment, and final inclusion. Identification and global deduplication are registered, but identification-set closure and final screening dispositions were not complete at this revision. The 140,959-record pool must not be presented as screened or eligible. Duplicate consolidation is a reconciliation step rather than an eligibility exclusion. Unresolved records and reports not retrieved require explicit accounting, separate from exclusions established by reviewed evidence.

Reproducibility materials comprise versioned protocols and amendments, executed queries, retrieval and import provenance, normalization rules, identity decisions, screening inputs and returns, and the derivation of reported counts. Stable identifiers connect the stages. Evidence enrichment and later adjudication retain their original inputs, enabling comparison before and after correction. The final availability statement must distinguish artifacts actually released from internal records and identify any redistribution restrictions. Code alone does not reproduce a changing provider index or a nondeterministic model response.

Coverage is limited by search terminology, indexed fields, provider coverage, unavailable metadata, and differences between API and local-snapshot matching. The retrieval cutoff is not a guarantee that every source was complete through that date. Prior-survey working materials and reference lists do not necessarily reconstruct the original surveys' included sets. English-language full-text requirements further restrict the population available for formal assessment. Broad searches and overlapping sources support candidate identification but do not establish exhaustive coverage or measured retrieval recall.

Screening introduces both interpretive and evidence-quality uncertainty. A model or human reviewer may infer an analytic relationship that is absent, overlook one that is poorly described, or assess mismatched provider metadata. The pilot's detected title/abstract conflict illustrates the need to distinguish source errors from reviewer disagreement. Preserved evidence and decisions make such errors inspectable but do not eliminate them. Findings from the mixed pilot sample and its non-independent re-analyses cannot establish corpus-wide accuracy or an unbiased estimate of inclusion prevalence.

\noindent\textbf{Completion required before submission:} Insert the reconciled review-flow diagram and final screening, unresolved, retrieval, and inclusion counts; report only validation analyses actually executed; and supply the final artifact-availability statement. Any remaining unresolved population must be disclosed when describing the scope and completeness of the synthesis.
